
\documentclass{article}
\usepackage{colortbl}
\usepackage{makecell}
\usepackage{multirow}
\usepackage{supertabular}

\begin{document}

\newcounter{utterance}

\centering \large Interaction Transcript for game `dond', experiment `semi\_en', episode 10 with Qwen3{-}30B{-}A3B{-}t0.0.
\vspace{24pt}

{ \footnotesize  \setcounter{utterance}{1}
\setlength{\tabcolsep}{0pt}
\begin{supertabular}{c@{$\;$}|p{.15\linewidth}@{}p{.15\linewidth}p{.15\linewidth}p{.15\linewidth}p{.15\linewidth}p{.15\linewidth}}
    \# & $\;$A & \multicolumn{4}{c}{Game Master} & $\;\:$B\\
    \hline

    \theutterance \stepcounter{utterance}  
    & & \multicolumn{4}{p{0.6\linewidth}}{
        \cellcolor[rgb]{0.9,0.9,0.9}{
            \makecell[{{p{\linewidth}}}]{
                \texttt{\tiny{[P1$\langle$GM]}}
                \texttt{You are playing a negotiation game in which you have to agree on how to divide a set of items among you and another player.} \\
\\ 
\texttt{Rules:} \\
\texttt{(a) You and the other player are given a set of items. Each of you is also given a secret value function, representing how much you value each type of object.} \\
\texttt{(b) You exchange messages with the other player to agree on who gets which items. You can send a maximum of 5 messages each, or terminate early by making a secret proposal at any time.} \\
\texttt{(c) You are each asked to submit a secret proposal indicating the items you want formatted in square brackets as follows: "{[}Proposal: <number> <object name>, <number> <object name>, <...>{]}"} \\
\texttt{(d) If your proposals are complementary, i.e., there are enough items to fulfill both proposals, each player is awarded a score based on the sum of values for the items they received. Otherwise, both of you get zero points.} \\
\texttt{(e) Your goal is to maximize the score you receive.} \\
\\ 
\texttt{Let us start.} \\
\\ 
\texttt{The set of available items is:} \\
\\ 
\texttt{1 bucket, 3 knives, 1 battery} \\
\\ 
\texttt{Your secret value function is:} \\
\\ 
\texttt{bucket: 2, knife: 2, battery: 2} \\
\\ 
\texttt{IMPORTANT: Your messages, unless it is the secret proposal, are directly transmitted to the other player, so do not include any response to the rules or text announcing your message. To make a secret proposal, use the indicated format. Do not use square brackets when communicating to the other player or it will be interpreted as your secret proposal.} \\
            }
        }
    }
    & \\ \\

    \theutterance \stepcounter{utterance}  
    & \multicolumn{4}{p{0.6\linewidth}}{
        \cellcolor[rgb]{0.9,0.9,0.9}{
            \makecell[{{p{\linewidth}}}]{
                \texttt{\tiny{[P1$\rangle$GM]}}
                \texttt{<think>} \\
\texttt{Okay, let's see. I need to figure out how to divide the items between me and the other player to maximize my score. The items are 1 bucket, 3 knives, and 1 battery. My value for each is 2, so each item is worth 2 points to me. The other player's values aren't given, so I have to guess or assume. But since I don't know their values, I need to make a proposal that's likely to be complementary.} \\
\\ 
\texttt{First, the rules say that if our proposals are complementary (enough items to fulfill both), we get points based on our own values. Otherwise, zero. So I need to make a proposal that the other player can accept without conflicting. But how?} \\
\\ 
\texttt{Since I don't know their values, maybe I should think about what's most valuable to me. Each item is 2 points, so I want to get as many as possible. But if I take all, the other player might not get anything, but maybe they can't. Wait, but the other player's proposal has to be complementary. So if I take all items, the other player's proposal would have to be empty, but that's not possible. So maybe I need to split the items in a way that both can get some.} \\
\\ 
\texttt{But how to communicate that? Since I can send up to 5 messages, but maybe I can just make a secret proposal directly. But the problem is that the other player's values are unknown. So maybe I should propose a fair split. For example, take 1 bucket, 1 knife, and 1 battery, and leave the rest. But wait, there's only 1 bucket and 1 battery. So maybe I can take the bucket and battery, and 1 knife, and the other player gets the remaining 2 knives. But that's just a guess.} \\
\\ 
\texttt{Alternatively, since each item is worth 2 to me, maybe I should take as many as possible. But if I take all, the other player can't get anything, so their proposal would have to be empty, which might not be allowed. But the rules don't specify that proposals have to have items. So maybe I can propose taking all items, and the other player can't do anything, but then their proposal would be empty, which might not be complementary. Wait, if I take all items, then the other player's proposal would have to be something that doesn't} \\
            }
        }
    }
    & & \\ \\

    \theutterance \stepcounter{utterance}  
    & & & \multicolumn{2}{p{0.3\linewidth}}{
        \cellcolor[rgb]{0.9,0.9,0.9}{
            \makecell[{{p{\linewidth}}}]{
                \texttt{\tiny{[GM$|$GM]}}
                \texttt{continue} \\
            }
        }
    }
    & & \\ \\

    \theutterance \stepcounter{utterance}  
    & & & \multicolumn{2}{p{0.3\linewidth}}{
        \cellcolor[rgb]{0.9,0.9,0.9}{
            \makecell[{{p{\linewidth}}}]{
                \texttt{\tiny{[GM$|$GM]}}
                \texttt{<think>} \\
\texttt{Okay, let's see. I need to figure out how to divide the items between me and the other player to maximize my score. The items are 1 bucket, 3 knives, and 1 battery. My value for each is 2, so each item is worth 2 points to me. The other player's values aren't given, so I have to guess or assume. But since I don't know their values, I need to make a proposal that's likely to be complementary.} \\
\\ 
\texttt{First, the rules say that if our proposals are complementary (enough items to fulfill both), we get points based on our own values. Otherwise, zero. So I need to make a proposal that the other player can accept without conflicting. But how?} \\
\\ 
\texttt{Since I don't know their values, maybe I should think about what's most valuable to me. Each item is 2 points, so I want to get as many as possible. But if I take all, the other player might not get anything, but maybe they can't. Wait, but the other player's proposal has to be complementary. So if I take all items, the other player's proposal would have to be empty, but that's not possible. So maybe I need to split the items in a way that both can get some.} \\
\\ 
\texttt{But how to communicate that? Since I can send up to 5 messages, but maybe I can just make a secret proposal directly. But the problem is that the other player's values are unknown. So maybe I should propose a fair split. For example, take 1 bucket, 1 knife, and 1 battery, and leave the rest. But wait, there's only 1 bucket and 1 battery. So maybe I can take the bucket and battery, and 1 knife, and the other player gets the remaining 2 knives. But that's just a guess.} \\
\\ 
\texttt{Alternatively, since each item is worth 2 to me, maybe I should take as many as possible. But if I take all, the other player can't get anything, so their proposal would have to be empty, which might not be allowed. But the rules don't specify that proposals have to have items. So maybe I can propose taking all items, and the other player can't do anything, but then their proposal would be empty, which might not be complementary. Wait, if I take all items, then the other player's proposal would have to be something that doesn't} \\
            }
        }
    }
    & & \\ \\

    \theutterance \stepcounter{utterance}  
    & & \multicolumn{4}{p{0.6\linewidth}}{
        \cellcolor[rgb]{0.9,0.9,0.9}{
            \makecell[{{p{\linewidth}}}]{
                \texttt{\tiny{[GM$\rangle$P2]}}
                \texttt{You are playing a negotiation game in which you have to agree on how to divide a set of items among you and another player.} \\
\\ 
\texttt{Rules:} \\
\texttt{(a) You and the other player are given a set of items. Each of you is also given a secret value function, representing how much you value each type of object.} \\
\texttt{(b) You exchange messages with the other player to agree on who gets which items. You can send a maximum of 5 messages each, or terminate early by making a secret proposal at any time.} \\
\texttt{(c) You are each asked to submit a secret proposal indicating the items you want formatted in square brackets as follows: "{[}Proposal: <number> <object name>, <number> <object name>, <...>{]}"} \\
\texttt{(d) If your proposals are complementary, i.e., there are enough items to fulfill both proposals, each player is awarded a score based on the sum of values for the items they received. Otherwise, both of you get zero points.} \\
\texttt{(e) Your goal is to maximize the score you receive.} \\
\\ 
\texttt{Let us start.} \\
\\ 
\texttt{The set of available items is:} \\
\\ 
\texttt{1 bucket, 3 knives, 1 battery} \\
\\ 
\texttt{Your secret value function is:} \\
\\ 
\texttt{bucket: 3, knife: 2, battery: 1} \\
\\ 
\texttt{IMPORTANT: Your messages, unless it is the secret proposal, are directly transmitted to the other player, so do not include any response to the rules or text announcing your message. To make a secret proposal, use the indicated format. Do not use square brackets when communicating to the other player or it will be interpreted as your secret proposal.} \\
\\ 
\\ 
\texttt{<think>} \\
\texttt{Okay, let's see. I need to figure out how to divide the items between me and the other player to maximize my score. The items are 1 bucket, 3 knives, and 1 battery. My value for each is 2, so each item is worth 2 points to me. The other player's values aren't given, so I have to guess or assume. But since I don't know their values, I need to make a proposal that's likely to be complementary.} \\
\\ 
\texttt{First, the rules say that if our proposals are complementary (enough items to fulfill both), we get points based on our own values. Otherwise, zero. So I need to make a proposal that the other player can accept without conflicting. But how?} \\
\\ 
\texttt{Since I don't know their values, maybe I should think about what's most valuable to me. Each item is 2 points, so I want to get as many as possible. But if I take all, the other player might not get anything, but maybe they can't. Wait, but the other player's proposal has to be complementary. So if I take all items, the other player's proposal would have to be empty, but that's not possible. So maybe I need to split the items in a way that both can get some.} \\
\\ 
\texttt{But how to communicate that? Since I can send up to 5 messages, but maybe I can just make a secret proposal directly. But the problem is that the other player's values are unknown. So maybe I should propose a fair split. For example, take 1 bucket, 1 knife, and 1 battery, and leave the rest. But wait, there's only 1 bucket and 1 battery. So maybe I can take the bucket and battery, and 1 knife, and the other player gets the remaining 2 knives. But that's just a guess.} \\
\\ 
\texttt{Alternatively, since each item is worth 2 to me, maybe I should take as many as possible. But if I take all, the other player can't get anything, so their proposal would have to be empty, which might not be allowed. But the rules don't specify that proposals have to have items. So maybe I can propose taking all items, and the other player can't do anything, but then their proposal would be empty, which might not be complementary. Wait, if I take all items, then the other player's proposal would have to be something that doesn't} \\
            }
        }
    }
    & \\ \\

    \theutterance \stepcounter{utterance}  
    & & & \multicolumn{4}{p{0.6\linewidth}}{
        \cellcolor[rgb]{0.9,0.9,0.9}{
            \makecell[{{p{\linewidth}}}]{
                \texttt{\tiny{[GM$\langle$P2]}}
                \texttt{<think>} \\
\texttt{Okay, let's try to think through this step by step. The goal is to maximize my score, which is the sum of the values of the items I get. The items are 1 bucket, 3 knives, 1 battery. My values are bucket: 3, knife: 2, battery: 1. So the bucket is the most valuable to me, then knives, then battery.} \\
\\ 
\texttt{But the other player's values are unknown. So I need to make a proposal that's likely to be complementary. If I take the bucket and some knives, maybe the other player would want the remaining knives or the battery. But how to communicate that?} \\
\\ 
\texttt{Since I can send messages, maybe I can start by proposing a split. But since I can only send 5 messages, maybe I should try to propose directly. However, if I make a secret proposal, it's possible that the other player's proposal might not complement mine. So I need to think about what's the most likely complementary proposal.} \\
\\ 
\texttt{Alternatively, maybe I should try to get the bucket and as many knives as possible. Let's say I propose taking the bucket and 2 knives. That would leave 1 knife and the battery. The other player might want the remaining knife and battery. But if they value the battery more, maybe they would take it. But without knowing their values, it's hard to tell.} \\
\\ 
\texttt{Wait, but the other player's proposal has to be complementary. So if I take the bucket and 2 knives, the other player can take the remaining 1 knife and battery. That would work. But how to ensure that? Maybe I should propose a split that leaves the other player with items that are valuable to them. But since I don't know their values, I have to make an assumption.} \\
\\ 
\texttt{Alternatively, maybe the other player's values are similar to mine. If they also value the bucket highly, then maybe they would want the bucket. But if I take the bucket, they might not get it. So maybe I should propose a split where I take the bucket and some knives, and they take the rest. But how to communicate that?} \\
\\ 
\texttt{Alternatively, maybe I should propose taking the bucket and the battery, and 1 knife. Then the other player can take the remaining 2 knives. That way, both of us get items. But again, without knowing their values, it's a guess.} \\
\\ 
\texttt{Another approach: since the total value of all items is 3} \\
            }
        }
    }
     \\ \\

    \theutterance \stepcounter{utterance}  
    & & & \multicolumn{2}{p{0.3\linewidth}}{
        \cellcolor[rgb]{0.9,0.9,0.9}{
            \makecell[{{p{\linewidth}}}]{
                \texttt{\tiny{[GM$|$GM]}}
                \texttt{continue} \\
            }
        }
    }
    & & \\ \\

    \theutterance \stepcounter{utterance}  
    & & & \multicolumn{2}{p{0.3\linewidth}}{
        \cellcolor[rgb]{0.9,0.9,0.9}{
            \makecell[{{p{\linewidth}}}]{
                \texttt{\tiny{[GM$|$GM]}}
                \texttt{<think>} \\
\texttt{Okay, let's try to think through this step by step. The goal is to maximize my score, which is the sum of the values of the items I get. The items are 1 bucket, 3 knives, 1 battery. My values are bucket: 3, knife: 2, battery: 1. So the bucket is the most valuable to me, then knives, then battery.} \\
\\ 
\texttt{But the other player's values are unknown. So I need to make a proposal that's likely to be complementary. If I take the bucket and some knives, maybe the other player would want the remaining knives or the battery. But how to communicate that?} \\
\\ 
\texttt{Since I can send messages, maybe I can start by proposing a split. But since I can only send 5 messages, maybe I should try to propose directly. However, if I make a secret proposal, it's possible that the other player's proposal might not complement mine. So I need to think about what's the most likely complementary proposal.} \\
\\ 
\texttt{Alternatively, maybe I should try to get the bucket and as many knives as possible. Let's say I propose taking the bucket and 2 knives. That would leave 1 knife and the battery. The other player might want the remaining knife and battery. But if they value the battery more, maybe they would take it. But without knowing their values, it's hard to tell.} \\
\\ 
\texttt{Wait, but the other player's proposal has to be complementary. So if I take the bucket and 2 knives, the other player can take the remaining 1 knife and battery. That would work. But how to ensure that? Maybe I should propose a split that leaves the other player with items that are valuable to them. But since I don't know their values, I have to make an assumption.} \\
\\ 
\texttt{Alternatively, maybe the other player's values are similar to mine. If they also value the bucket highly, then maybe they would want the bucket. But if I take the bucket, they might not get it. So maybe I should propose a split where I take the bucket and some knives, and they take the rest. But how to communicate that?} \\
\\ 
\texttt{Alternatively, maybe I should propose taking the bucket and the battery, and 1 knife. Then the other player can take the remaining 2 knives. That way, both of us get items. But again, without knowing their values, it's a guess.} \\
\\ 
\texttt{Another approach: since the total value of all items is 3} \\
            }
        }
    }
    & & \\ \\

    \theutterance \stepcounter{utterance}  
    & & \multicolumn{4}{p{0.6\linewidth}}{
        \cellcolor[rgb]{0.9,0.9,0.9}{
            \makecell[{{p{\linewidth}}}]{
                \texttt{\tiny{[P1$\langle$GM]}}
                \texttt{<think>} \\
\texttt{Okay, let's try to think through this step by step. The goal is to maximize my score, which is the sum of the values of the items I get. The items are 1 bucket, 3 knives, 1 battery. My values are bucket: 3, knife: 2, battery: 1. So the bucket is the most valuable to me, then knives, then battery.} \\
\\ 
\texttt{But the other player's values are unknown. So I need to make a proposal that's likely to be complementary. If I take the bucket and some knives, maybe the other player would want the remaining knives or the battery. But how to communicate that?} \\
\\ 
\texttt{Since I can send messages, maybe I can start by proposing a split. But since I can only send 5 messages, maybe I should try to propose directly. However, if I make a secret proposal, it's possible that the other player's proposal might not complement mine. So I need to think about what's the most likely complementary proposal.} \\
\\ 
\texttt{Alternatively, maybe I should try to get the bucket and as many knives as possible. Let's say I propose taking the bucket and 2 knives. That would leave 1 knife and the battery. The other player might want the remaining knife and battery. But if they value the battery more, maybe they would take it. But without knowing their values, it's hard to tell.} \\
\\ 
\texttt{Wait, but the other player's proposal has to be complementary. So if I take the bucket and 2 knives, the other player can take the remaining 1 knife and battery. That would work. But how to ensure that? Maybe I should propose a split that leaves the other player with items that are valuable to them. But since I don't know their values, I have to make an assumption.} \\
\\ 
\texttt{Alternatively, maybe the other player's values are similar to mine. If they also value the bucket highly, then maybe they would want the bucket. But if I take the bucket, they might not get it. So maybe I should propose a split where I take the bucket and some knives, and they take the rest. But how to communicate that?} \\
\\ 
\texttt{Alternatively, maybe I should propose taking the bucket and the battery, and 1 knife. Then the other player can take the remaining 2 knives. That way, both of us get items. But again, without knowing their values, it's a guess.} \\
\\ 
\texttt{Another approach: since the total value of all items is 3} \\
            }
        }
    }
    & \\ \\

    \theutterance \stepcounter{utterance}  
    & \multicolumn{4}{p{0.6\linewidth}}{
        \cellcolor[rgb]{0.9,0.9,0.9}{
            \makecell[{{p{\linewidth}}}]{
                \texttt{\tiny{[P1$\rangle$GM]}}
                \texttt{<think>} \\
\texttt{Okay, let's try to figure this out. The items are 1 bucket, 3 knives, 1 battery. My values are bucket: 3, knife: 2, battery: 1. The other player's values are unknown, but I need to make a proposal that's complementary.} \\
\\ 
\texttt{If I take the bucket (3 points) and 2 knives (2*2=4), that's 7 points. The remaining items are 1 knife and 1 battery. If the other player takes those, they get 2 + 1 = 3 points. But maybe they value the battery more? Or maybe they want more knives. Alternatively, if I take the bucket and 1 knife (3+2=5), leaving 2 knives and 1 battery. The other player could take those, getting 2*2 +1=5. That's a balanced split.} \\
\\ 
\texttt{But how to ensure that the other player's proposal complements mine. Since I don't know their values, maybe I should propose a split that leaves them with items that are valuable to them. But without knowing their values, it's tricky. Maybe the safest is to take the bucket and 1 knife, and let them have the rest. That way, they get 2 knives and a battery, which might be 2*2 +1=5. But if they value the battery more, maybe they'd prefer that. Alternatively, if they take the battery and 1 knife, that's 1+2=3. But I don't know.} \\
\\ 
\texttt{Alternatively, maybe I should propose taking the bucket and the battery, which gives me 3+1=4, and leave 3 knives. The other player could take all 3 knives, getting 3*2=6. That's a possible split. But if they value the bucket more, they might want it, but I already took it. So maybe that's not complementary.} \\
\\ 
\texttt{Wait, the key is that the proposals need to be complementary. So if I propose taking the bucket and 2 knives, and the other player proposes taking the remaining 1 knife and battery, that works. But how to ensure that? Since I can't communicate their values, maybe I should just make a proposal that's likely to be complementary. Maybe the most straightforward is to take the bucket and 1 knife, and let them have the rest. That way, they get 2 knives and a} \\
            }
        }
    }
    & & \\ \\

    \theutterance \stepcounter{utterance}  
    & & & \multicolumn{2}{p{0.3\linewidth}}{
        \cellcolor[rgb]{0.9,0.9,0.9}{
            \makecell[{{p{\linewidth}}}]{
                \texttt{\tiny{[GM$|$GM]}}
                \texttt{continue} \\
            }
        }
    }
    & & \\ \\

    \theutterance \stepcounter{utterance}  
    & & & \multicolumn{2}{p{0.3\linewidth}}{
        \cellcolor[rgb]{0.9,0.9,0.9}{
            \makecell[{{p{\linewidth}}}]{
                \texttt{\tiny{[GM$|$GM]}}
                \texttt{<think>} \\
\texttt{Okay, let's try to figure this out. The items are 1 bucket, 3 knives, 1 battery. My values are bucket: 3, knife: 2, battery: 1. The other player's values are unknown, but I need to make a proposal that's complementary.} \\
\\ 
\texttt{If I take the bucket (3 points) and 2 knives (2*2=4), that's 7 points. The remaining items are 1 knife and 1 battery. If the other player takes those, they get 2 + 1 = 3 points. But maybe they value the battery more? Or maybe they want more knives. Alternatively, if I take the bucket and 1 knife (3+2=5), leaving 2 knives and 1 battery. The other player could take those, getting 2*2 +1=5. That's a balanced split.} \\
\\ 
\texttt{But how to ensure that the other player's proposal complements mine. Since I don't know their values, maybe I should propose a split that leaves them with items that are valuable to them. But without knowing their values, it's tricky. Maybe the safest is to take the bucket and 1 knife, and let them have the rest. That way, they get 2 knives and a battery, which might be 2*2 +1=5. But if they value the battery more, maybe they'd prefer that. Alternatively, if they take the battery and 1 knife, that's 1+2=3. But I don't know.} \\
\\ 
\texttt{Alternatively, maybe I should propose taking the bucket and the battery, which gives me 3+1=4, and leave 3 knives. The other player could take all 3 knives, getting 3*2=6. That's a possible split. But if they value the bucket more, they might want it, but I already took it. So maybe that's not complementary.} \\
\\ 
\texttt{Wait, the key is that the proposals need to be complementary. So if I propose taking the bucket and 2 knives, and the other player proposes taking the remaining 1 knife and battery, that works. But how to ensure that? Since I can't communicate their values, maybe I should just make a proposal that's likely to be complementary. Maybe the most straightforward is to take the bucket and 1 knife, and let them have the rest. That way, they get 2 knives and a} \\
            }
        }
    }
    & & \\ \\

    \theutterance \stepcounter{utterance}  
    & & \multicolumn{4}{p{0.6\linewidth}}{
        \cellcolor[rgb]{0.9,0.9,0.9}{
            \makecell[{{p{\linewidth}}}]{
                \texttt{\tiny{[GM$\rangle$P2]}}
                \texttt{<think>} \\
\texttt{Okay, let's try to figure this out. The items are 1 bucket, 3 knives, 1 battery. My values are bucket: 3, knife: 2, battery: 1. The other player's values are unknown, but I need to make a proposal that's complementary.} \\
\\ 
\texttt{If I take the bucket (3 points) and 2 knives (2*2=4), that's 7 points. The remaining items are 1 knife and 1 battery. If the other player takes those, they get 2 + 1 = 3 points. But maybe they value the battery more? Or maybe they want more knives. Alternatively, if I take the bucket and 1 knife (3+2=5), leaving 2 knives and 1 battery. The other player could take those, getting 2*2 +1=5. That's a balanced split.} \\
\\ 
\texttt{But how to ensure that the other player's proposal complements mine. Since I don't know their values, maybe I should propose a split that leaves them with items that are valuable to them. But without knowing their values, it's tricky. Maybe the safest is to take the bucket and 1 knife, and let them have the rest. That way, they get 2 knives and a battery, which might be 2*2 +1=5. But if they value the battery more, maybe they'd prefer that. Alternatively, if they take the battery and 1 knife, that's 1+2=3. But I don't know.} \\
\\ 
\texttt{Alternatively, maybe I should propose taking the bucket and the battery, which gives me 3+1=4, and leave 3 knives. The other player could take all 3 knives, getting 3*2=6. That's a possible split. But if they value the bucket more, they might want it, but I already took it. So maybe that's not complementary.} \\
\\ 
\texttt{Wait, the key is that the proposals need to be complementary. So if I propose taking the bucket and 2 knives, and the other player proposes taking the remaining 1 knife and battery, that works. But how to ensure that? Since I can't communicate their values, maybe I should just make a proposal that's likely to be complementary. Maybe the most straightforward is to take the bucket and 1 knife, and let them have the rest. That way, they get 2 knives and a} \\
            }
        }
    }
    & \\ \\

    \theutterance \stepcounter{utterance}  
    & & & \multicolumn{4}{p{0.6\linewidth}}{
        \cellcolor[rgb]{0.9,0.9,0.9}{
            \makecell[{{p{\linewidth}}}]{
                \texttt{\tiny{[GM$\langle$P2]}}
                \texttt{<think>} \\
\texttt{Okay, let's try to make a plan. The items are 1 bucket, 3 knives, 1 battery. My values: bucket 3, knife 2, battery 1. The other player's values are unknown, but I need to propose something that's complementary.} \\
\\ 
\texttt{If I take the bucket (3 points) and 1 knife (2 points), that's 5 points. The remaining items are 2 knives and 1 battery. The other player could take those, getting 2*2 +1 = 5 points. That's a balanced split. But maybe they want more knives? Or maybe they value the battery more. Alternatively, if I take the bucket and 2 knives (3+4=7), leaving 1 knife and battery (2+1=3). But if the other player takes those, they get 3 points. That might not be as good for them, but maybe they have to accept it.} \\
\\ 
\texttt{Alternatively, if I take the bucket and battery (3+1=4), leaving 3 knives (6 points). The other player could take all 3 knives, getting 6 points. That's a possible split. But if they value the bucket more, they might not want that. But since I don't know, maybe this is a safe bet.} \\
\\ 
\texttt{But how to decide? Since the other player's values are unknown, maybe the best approach is to take the most valuable item (bucket) and then some knives, leaving the rest for them. Let's say I take the bucket and 1 knife (total 5), and they take 2 knives and battery (5). That way, both get 5 points. But if they value the battery more, maybe they'd prefer that. Alternatively, if I take the bucket and 2 knives (7), they get 1 knife and battery (3). But maybe they can't get more than that.} \\
\\ 
\texttt{Alternatively, maybe the other player would prefer to take the battery and some knives. But without knowing their values, it's hard. The safest is to propose a split where I take the bucket and 1 knife, and they get the rest. That way, the items are divided, and both can get something. So my secret proposal would be {[}Proposal: 1 bucket, 1 knife{]}. Then the other player's proposal would need to take the remaining 2 knives and 1 battery. If that's possible, then both} \\
            }
        }
    }
     \\ \\

    \theutterance \stepcounter{utterance}  
    & & & \multicolumn{2}{p{0.3\linewidth}}{
        \cellcolor[rgb]{0.9,0.9,0.9}{
            \makecell[{{p{\linewidth}}}]{
                \texttt{\tiny{[GM$|$GM]}}
                \texttt{proposal} \\
            }
        }
    }
    & & \\ \\

    \theutterance \stepcounter{utterance}  
    & & & \multicolumn{2}{p{0.3\linewidth}}{
        \cellcolor[rgb]{0.9,0.9,0.9}{
            \makecell[{{p{\linewidth}}}]{
                \texttt{\tiny{[GM$|$GM]}}
                [1, 1, 0]
            }
        }
    }
    & & \\ \\

    \theutterance \stepcounter{utterance}  
    & & \multicolumn{4}{p{0.6\linewidth}}{
        \cellcolor[rgb]{0.9,0.9,0.9}{
            \makecell[{{p{\linewidth}}}]{
                \texttt{\tiny{[P1$\langle$GM]}}
                \texttt{The other player has submitted a secret proposal. It is now your turn to submit a proposal of your own.} \\
            }
        }
    }
    & \\ \\

    \theutterance \stepcounter{utterance}  
    & \multicolumn{4}{p{0.6\linewidth}}{
        \cellcolor[rgb]{0.9,0.9,0.9}{
            \makecell[{{p{\linewidth}}}]{
                \texttt{\tiny{[P1$\rangle$GM]}}
                \texttt{<think>} \\
\texttt{Okay, the other player has already submitted their secret proposal, and now it's my turn. I need to figure out what their proposal might be and then make a complementary one. But since I don't know their values, I have to guess.} \\
\\ 
\texttt{First, let's recall the items: 1 bucket, 3 knives, 1 battery. My values are bucket: 3, knife: 2, battery: 1. The other player's values are unknown, but maybe they have similar or different preferences.} \\
\\ 
\texttt{If I assume that the other player might also value the bucket highly, they might have proposed taking the bucket. Alternatively, they might have gone for the knives since there are more of them. Or maybe the battery, but it's only worth 1 to me.} \\
\\ 
\texttt{Since I can't know for sure, I need to make a proposal that can complement whatever they proposed. Let's think about possible scenarios.} \\
\\ 
\texttt{Case 1: The other player proposed taking the bucket. Then I should not take the bucket. So I could take knives and battery. But if they took the bucket, I need to make sure I don't conflict. So maybe propose taking 2 knives and the battery. That would leave 1 knife for them, but if they took the bucket, maybe they don't need the knife. But this is speculative.} \\
\\ 
\texttt{Case 2: The other player proposed taking all knives. Then I should take the bucket and battery. That way, we don't conflict. But if they took all knives, I can take the bucket and battery, which are 3 +1 =4 points. That's better than nothing.} \\
\\ 
\texttt{Case 3: The other player proposed taking the battery. Then I can take the bucket and knives. But again, not sure.} \\
\\ 
\texttt{Alternatively, maybe the other player's proposal is something like {[}1 bucket, 1 knife{]}, and I need to complement that by taking the remaining 2 knives and the battery. But how to know?} \\
\\ 
\texttt{Since I have to make a proposal now, I need to choose a combination that's likely to be complementary. Let's think about what's the most valuable for me. The bucket is 3, which is more than a knife (2) or battery (1). So if I can get the bucket, that's good. But if the other player also wants the bucket, we can't both have it. So maybe I should propose taking the bucket and some knives, assuming they don't} \\
            }
        }
    }
    & & \\ \\

    \theutterance \stepcounter{utterance}  
    & & & \multicolumn{2}{p{0.3\linewidth}}{
        \cellcolor[rgb]{0.9,0.9,0.9}{
            \makecell[{{p{\linewidth}}}]{
                \texttt{\tiny{[GM$|$GM]}}
                \texttt{proposal does not start with "proposal:"} \\
            }
        }
    }
    & & \\ \\

\end{supertabular}
}

\end{document}
